\documentclass{article}
\usepackage{graphicx} % Required for inserting images
\usepackage{amssymb}
\usepackage{amsmath}

\newcommand{\bv}{\mathbf{v}}
\newcommand{\bu}{\mathbf{u}}
\newcommand{\bw}{\mathbf{w}}
\newcommand{\bO}{\mathbf{0}}
\newcommand{\R}{\mathbb{R}}
\newcommand{\Q}{\mathbb{Q}}
\newcommand{\C}{\mathbb{C}}
\newcommand{\F}{\mathbb{F}}

\title{Fundamentals Chapter 7}
\author{Siddharth K}
\date{November 2024}

\begin{document}

\maketitle

\textbf{1a.} $$\begin{bmatrix}
    0 & -1 \\
    1 & 0
\end{bmatrix}
\begin{bmatrix}
    0 & -1 \\
    1 & 0
\end{bmatrix} = 
\begin{bmatrix}
    -1 & 0 \\
    0 & -1
\end{bmatrix}$$

\textbf{1b.} $$\begin{bmatrix}
    2 & 3 & 1 \\
    0 & 4 & 1
\end{bmatrix}
\begin{bmatrix}
    2 & 2 \\
    1 & 1 \\
    1 & 0
\end{bmatrix} = 
\begin{bmatrix}
    8 & 7 \\
    5 & 4
\end{bmatrix}$$

\textbf{1c.} $$\begin{bmatrix}
    1 & 2 & 0 \\
    3 & 2 & 1 \\
    2 & 1 & 1
\end{bmatrix}
\begin{bmatrix}
    1 & 3 & 2 \\
    2 & 2 & 1 \\
    0 & 1 & 1
\end{bmatrix} = 
\begin{bmatrix}
    5 & 7 & 4 \\
    7 & 14 & 9 \\
    4 & 9 & 6
\end{bmatrix}$$

\textbf{1d.} First, we see that:
$$\begin{bmatrix}
    1 & 2 & 0 \\
    3 & 2 & 1 \\
    2 & 1 & 1
\end{bmatrix}\left(
\begin{bmatrix}
    1 & 3 & 2 \\
    2 & 2 & 1 \\
    0 & 1 & 1
\end{bmatrix}
\begin{bmatrix}
    x \\
    y \\
    z
\end{bmatrix}\right)$$
$$= \begin{bmatrix}
    1 & 2 & 0 \\
    3 & 2 & 1 \\
    2 & 1 & 1
\end{bmatrix}
\begin{bmatrix}
    x + 3y + 2z \\
    2x + 2y + z \\
    y + z
\end{bmatrix}
= \begin{bmatrix}
    5x + 7y + 4z \\
    7x + 14y + 9z \\
    4x + 9y + 6z
\end{bmatrix}.$$
Then:
$$\left(\begin{bmatrix}
    1 & 2 & 0 \\
    3 & 2 & 1 \\
    2 & 1 & 1
\end{bmatrix}
\begin{bmatrix}
    1 & 3 & 2 \\
    2 & 2 & 1 \\
    0 & 1 & 1
\end{bmatrix}\right)
\begin{bmatrix}
    x \\
    y \\
    z
\end{bmatrix}
= \begin{bmatrix}
    5 & 7 & 4 \\
    7 & 14 & 9 \\
    4 & 9 & 6
\end{bmatrix}
\begin{bmatrix}
    x \\
    y \\
    z
\end{bmatrix}$$
$$= \begin{bmatrix}
    5x + 7y + 4z \\
    7x + 14y + 9z \\
    4x + 9y + 6z
\end{bmatrix},$$
which matches the first one. \\

\textbf{2.} The transformation represented by $R_\theta$ is a rotation counterclockwise by $\theta$ radians. This is because carrying out the multiplication
$$\begin{bmatrix}
    \cos\theta & -\sin\theta \\
    \sin\theta & \cos\theta
\end{bmatrix}
\begin{bmatrix}
    x \\
    y
\end{bmatrix}$$
gives us
$$\begin{bmatrix}
    x\cos\theta - y\sin\theta \\
    x\sin\theta + y\cos\theta
\end{bmatrix}$$
which is the formula for rotation of the point $(x, y)$ about the origin by $\theta$ radians. \\

\textbf{3a.} We can uniquely write $\bv = a_1\bv_1 + \cdots + a_n\bv_n$. So thinking of $\bv$ as $(a_1, \ldots, a_n)$ means uniquely representing $\bv$ by its coordinates in $\F^n$ with respect to the given basis. \\

\textbf{3b.} Assume that $(a_1, \ldots, a_n)_V = (b_1, \ldots, b_n)_W$ for all $a_i, b_i \in \F$. Then any vector $\bv \in \F^n$ has the same coordinates whether it's represented in the basis $B_V$ or $B_U$. Now, let $(a_1, \ldots, a_n) = (b_1, \ldots b_n) = (0, \ldots, 0, 1, 0 \ldots, 0)$ where $a_i = b_i = 1$. By the assumption, we have that $\bv_i = \bw_i$. Doing this for all $i$ gives us $\bv_i = \bw_i$ for all $i$. \\

Now, assume that $\bv_i = \bw_i$. There is only one way to write a vector $\bv$ in terms of $\bv_i$, so $(a_1, \ldots, a_n)_V = (b_1, \ldots, b_n)_W$. \\

\textbf{4.} Since $T(\bu_i) = \bw_i$, given any $\bw \in W$, we have that
$$\bw = \sum a_i\bw_i = \sum a_iT(\bu_i) = T\left(\sum a_i\bu_i\right) = T(\bu)$$
for some $\bu \in U$. So $T$ is invertible. \\

\textbf{5a.} By the definition of matrix multiplication, $C = AB$ is defined as
$$c_{ij} = \sum_{l = 1}^n a_{il}b_{lj}.$$
Since $i$ and $j$ range from $1$ to $n$, there are $n$ values for $i$ and $j$ in $c_{ij}$, so $C$ is an $n\times n$ matrix. \\

\textbf{5b.} We have that the entry at $(i, j)$ in $AI_n$ is the total sum of each entry in
$\begin{bmatrix}
    a_{i1} & \ldots & a_{in}
\end{bmatrix}$
multiplied by each entry in
$\begin{bmatrix}
    0 & \ldots & 1 & \ldots & 0
\end{bmatrix}$ (where the $1$ is at position $j$), which is just $a_{ij}$. So $A = AI_n$. \\

Also, we have that the entry at $(i, j)$ in $I_nA$ is the total sum of each entry in
$\begin{bmatrix}
    0 & \ldots & 1 & \ldots & 0
\end{bmatrix}$
multiplied by each entry in
$\begin{bmatrix}
    a_{1j} & \ldots & a_{nj}
\end{bmatrix}$ (where the $1$ is at position $i$), which is just $a_{ij}$. So $A = I_nA$. \\

\textbf{5c.} First, consider the product $A(BC)$. We have that $$A(BC)_{ij} = \sum_{k = 1}^p a_{ik}(BC)_{kj} = \sum_{k = 1}^p a_{ik}\left(\sum_{l = 1}^n b_{kl}c_{lj}\right)$$
$$= \sum_{l = 1}^n\left(\sum_{k = 1}^p a_{ik}b_{kl})c_{lj}\right).$$
Now, consider the product $(AB)C$. We have that
$$(AB)C_{ij} = \sum_{l = 1}^n\left(\sum_{k = 1}^p a_{ik}b_{kl})c_{lj}\right),$$
which matches the expression for $A(BC)_{ij}$. Since this is true for all $i, j$, we have that $A(BC) = (AB)C$.\\

\textbf{5d.} Let $A = \begin{bmatrix} a & b \\ c & d \end{bmatrix}$ and $B = \begin{bmatrix} e & f \\ g & h \end{bmatrix}$. The condition is that $(a, b)$ is not a scalar multiple of $(c, d)$. This is because we need the system of equations created by $AB = I_2$ to have a unique solution, which is only true when $(a, b)$ is not a scalar multiple of $(c, d)$. \\

\textbf{5e.} Any matrix $A$ that isn't a scalar multiple of $I$ doesn't guarantee that $AB = BA$ for all $2\times 2$ matrices $B$. To see this, we can consider a few cases for $B$. Let $B$ equal
$\begin{bmatrix}
    1 & 0 \\
    0 & 0
\end{bmatrix}.$
Then
$$AB = \begin{bmatrix}
    a & b \\
    c & d
\end{bmatrix}
\begin{bmatrix}
    1 & 0 \\
    0 & 0
\end{bmatrix}
= \begin{bmatrix}
    a & 0 \\
    c & 0
\end{bmatrix},$$
and 
$$BA = \begin{bmatrix}
    1 & 0 \\
    0 & 0
\end{bmatrix}
\begin{bmatrix}
    a & b \\
    c & d
\end{bmatrix}
= \begin{bmatrix}
    a & b \\
    0 & 0
\end{bmatrix},$$
so we have that $b = c = 0$. This tells us that $A$ is of the form
$\begin{bmatrix}
    a & 0 \\
    0 & d
\end{bmatrix}.$
Now, letting $B$ equal
$\begin{bmatrix}
    0 & 1 \\
    0 & 0
\end{bmatrix}$
gives us
$$AB = \begin{bmatrix}
    a & 0 \\
    0 & d
\end{bmatrix}
\begin{bmatrix}
    0 & 1 \\
    0 & 0
\end{bmatrix}
= \begin{bmatrix}
    0 & a \\
    0 & 0
\end{bmatrix},$$
and
$$BA = \begin{bmatrix}
    0 & 1 \\
    0 & 0
\end{bmatrix}
\begin{bmatrix}
    a & 0 \\
    0 & d
\end{bmatrix}
= \begin{bmatrix}
    0 & d \\
    0 & 0
\end{bmatrix},$$
so $a = d$. This tells us that $A$ is of the form
$\begin{bmatrix}
    a & 0 \\
    0 & a
\end{bmatrix}.$
Furthermore, multiplying this by any $2\times 2$ matrix $B$ reveals that $AB = BA$. \\

\textbf{6.} The system of equations can be written as
$$A\bv = \begin{bmatrix}
    3 & -5 & 2 \\
    -1 & 2 & 2 \\
    2 & -1 & -1
\end{bmatrix}
\begin{bmatrix}
    x \\
    y \\
    z
\end{bmatrix}
= \begin{bmatrix}
    1 \\
    -1 \\
    0
\end{bmatrix}.$$
If you know the inverse of $A$, you could multiply both sides by $A^{-1}$, and the result would be
$$\begin{bmatrix}
    x \\
    y \\
    z
\end{bmatrix}
= \begin{bmatrix}
    3 & -5 & 2 \\
    -1 & 2 & 2 \\
    2 & -1 & -1
\end{bmatrix}^{-1}
\begin{bmatrix}
    1 \\
    -1 \\
    0
\end{bmatrix}.$$

\textbf{7.} Lets say $U = \{\bu_i\}, W = \{\bw_i\}$ are bases for $V$. The matrix $P_U^W$ maps each $\bu_i$ to $\bw_i$, so its inverse maps each $\bw_i$ to $\bu_i$. But this is exactly $P_W^U$, so $(P_U^W)^{-1}$ and $P_W^U$ map each basis vector in $W$ to the same one in $U$, which means they are the same transformation. \\

\textbf{8.} Reflexivity: $A = IAI$ so $A\sim A$. \\

Symmetry: $A \sim B$ implies that $A = PBP^{-1}$ for some matrix $P$. Multiplying both sides on the left by $P^{-1}$ and on the right by $P$ gives us $B = P^{-1}AP$, so $B \sim A$. \\

Transitivity: $A \sim B, B\sim C$ implies that $A = PBP^{-1}$ and $B = QCQ^{-1}$ for some matrices $P, Q$. Substituting the expression for $B$ into $A$, we get $A = PQCQ^{-1}P^{-1}$. Since $(PQ)^{-1} = Q^{-1}P^{-1}$, we have that $A \sim C$. \\

\textbf{9.} Let $I_n = PAP^{-1}$. Multiplying both sides by $P^{-1}$ on the left and $P$ on the right gives us $A = P^{-1}P = I_n$. So $A = I_n$ is the only solution to $I_n = PAP^{-1}$. \\

\textbf{10.} We have that
$$\begin{bmatrix}
    0 & 1/2 \\
    2 & 0
\end{bmatrix}^2
= \begin{bmatrix}
    1 & 0 \\
    0 & 1
\end{bmatrix} = I,$$
so $A^{10} = I$. \\

\textbf{11.} Consider $A\bv = \bw$ where the $k$'th column of $A$ is all zeroes. Expanding the product shows that the result $\bw$ is calculated in terms of every entry of $\bv$ except $v_k$. Which means, if we wanted to map $\bw$ to some unique $\bv$, we couldn't because $v_k$ can vary freely but $A\bv$ will still equal $\bw$. So $A$ is not invertible. \\

Now, consider $A\bv = \bw$ where the $k$'th row of $A$ is all zeroes. Since the $k$'th entry of $\bw$ is the dot product of the $k$'th row and $\bv$, it will always be $0$. So any vector in $W$ with a $k$'th entry other than $0$ isn't reached by the linear transformation, and $A$ is not invertible. \\

\textbf{12.} If $\ker(A)$ is nontrivial then $A$ is not invertible. This is because $0$ can then be reverse-mapped to more than one vector, which is a contradiction. \\

Let's say $\bv_i$ is the basis that is being mapped to $\bw_i$ (the columns of $A$) by the linear transformation, which we can call $T$. Since $\{\bw_i\}$ is not linearly independent, there exist nonzero $a_i$ such that
$$\sum a_i\bw_i = 0.$$
Since $\bw_i = T(\bv_i)$,
$$\sum a_i T(\bv_i) = T\left(\sum a_i \bv_i\right) = 0.$$
Since $a_i \neq 0$ and $\bv_i$ are linearly independent, $\sum a_i \bv_i \neq 0$. So $T(0) = 0$ and $T\left(\sum a_i \bv_i\right) = 0$, which means the kernel is nontrivial and $A$ is not invertible.

\end{document}